\documentclass[10pt,a4paper]{amsart}
\usepackage[margin=19mm]{geometry}
\usepackage{amsmath,amssymb,mathtools}
\usepackage[scr=rsfs,cal=euler]{mathalfa}
\usepackage{xfrac}
\usepackage[unicode,hidelinks,bookmarks=false]{hyperref}
\newcommand{\ie}{i.e.\ }
\newcommand{\cf}{cf.\ }
\newcommand{\resp}{respectively\ }
\newcommand{\Subsection}{\subsection}
\providecommand{\nobreakdash}{\nobreak\hbox{-}}
\setlength{\emergencystretch}{2em}
\multlinegap0pt
\usepackage{bm}
\usepackage{bbm}
\usepackage{dsfont}
\usepackage{xspace}
\usepackage{cancel}
\usepackage{mathtools}
\usepackage{amsthm}
\makeatletter
\@ifundefined{claim}{\newtheorem{claim}{Claim}}{}
\@ifundefined{lem}{\newtheorem{lem}[claim]{Lemma}}{}
\@ifundefined{thm}{\newtheorem{thm}{Theorem}}{}
\makeatother
\newcommand{\C}{{\mathbb C}}
\newcommand{\Gal}{{\mathrm{Gal}}}
\newcommand{\Qb}{{\overline{\Q}}}
\newcommand{\Q}{{\mathbb Q}}
\newcommand{\Z}{{\mathbb Z}}
\newcommand{\case}[1]{\vspace{0.3cm}\noindent\framebox{#1.}}
\newcommand{\disc}{{\mathrm{disc}}}
\renewcommand{\Cup}{\bigcup}
\title[Lower Bounds on the Least Common Multiple of a Polynomial Se]{Lower Bounds on the Least Common Multiple of a Polynomial Sequence and its Radical}
\author{Alexei Entin}
\date{Source: arXiv:2401.05184v4 (CC BY 4.0)}
\begin{document}
\maketitle

\noindent\emph{Source and licence.} Lower Bounds on the Least Common Multiple of a Polynomial Sequence and its Radical. Version arXiv:2401.05184v4 (CC BY 4.0), distributed under the Creative Commons Attribution 4.0 International licence, as recorded in the arXiv licence metadata for this exact version and shown on the abstract page https://arxiv.org/abs/2401.05184v4. Full rights notice: licenses/arxiv-2401.05184-CC-BY-4.0.txt.\par\smallskip

\noindent\textit{Editorial scope.} Counted unit: the Lem whose source label is \texttt{lem: pre lem main2} in the article (source0.tex lines 416--419, source label \texttt{lem: pre lem main2}), delivered together with the article's own complete original proof of it (source0.tex lines 421--456, one \texttt{proof} environment of 36 lines). No mathematics is written, repaired or re-derived: statement and proof are the article's own text line for line. Required context retained uncounted: thm: main2 (lines 169--173); lem: hom (lines 308--313); lem: main1 (lines 325--326); lem: int comb (lines 409--409). Every definition, display and label of those lines is retained unchanged. Omitted from this package and not redistributed: 1--168, 174--307, 314--324, 327--408, 410--415, 420--420, 457--477. Nothing inside a retained excerpt is omitted or altered: every retained body is delivered exactly as the article's lines are written, so no omission receipt of any kind accompanies this package. Citations: the retained excerpts carry no citation command at all, so no bibliography entry of the article is transcribed and no reference is redistributed. Reference adaptation: none was needed and none was made; every reference inside the delivered unit resolves inside it, so no unresolved reference or citation is produced. Printed slips kept literal: no printed word, symbol, hypothesis or number of the counted statement or of its proof was altered. Layout and class: the article's own class is \texttt{amsart} with packages including enumerate, longtable, amsmath, amscd, amsfonts, amsthm, amssymb, latexsym, comment, stmaryrd, graphicx, dsfont, amsaddr, mathtools; the wrapper is the lane's pinned \texttt{amsart} editorial wrapper at 10pt on A4, which restates the theorem-like environments the retained text needs and adds the article's own macro definitions that the retained text needs (\texttt{\textbackslash C}, \texttt{\textbackslash Cup}, \texttt{\textbackslash Gal}, \texttt{\textbackslash Q}, \texttt{\textbackslash Qb}, \texttt{\textbackslash Z}, \texttt{\textbackslash case}, \texttt{\textbackslash disc}), with extra wrapper packages (bm, bbm, dsfont, xspace, cancel, mathtools). The delivered numbering is the unit's own. Assets: no retained span contains a figure environment, a graphics-inclusion command, a PDF-page inclusion, a table or a TikZ picture; no image file is copied into this package, the package declares no asset dependency and no image file is redistributed with it. Licence: version arXiv:2401.05184v4 of this article is distributed under the Creative Commons Attribution 4.0 International licence, as recorded in the arXiv licence metadata for this exact version and shown on the abstract page https://arxiv.org/abs/2401.05184v4. A full rights notice is delivered at licenses/arxiv-2401.05184-CC-BY-4.0.txt. Characters: every retained excerpt is pure ASCII, so the delivered engine, XeLaTeX with its default Latin Modern text font, reports no missing character.
\begin{thm}\label{thm: main2}
Let $f\in\Z[x],\deg f=d\ge 3$ be irreducible, $S\subset\overline\Q$ its set of roots, $G_f$ its Galois group, $e=|G_f|$. Suppose that one can partition $S$ into (non-empty) potent subsets $S_1,\ldots,S_r$ where $r$ satisfies $$\dim\langle S\cup\{1\}\rangle_\Q=d+1-r,$$ here $\langle\cdot\rangle_\Q$ denotes the linear span over $\Q$ (note that since the $S_i$ are potent and disjoint, each contributes a new linear dependency on $S\cup\{1\}$ and a priori $\dim\langle S\cup\{1\}\rangle_\Q\le d+1-r$).
Then $$\log\ell_f(N)\gtrsim \frac {d-1}{U_{d,e}}N\log N,$$
where \begin{equation}\label{eq: def U_e}U_{d,e}=\sum_{k=1}^{\min(d-2,\lfloor 3d/4\rfloor)}\min \left(d-k,\bigg\lceil\frac {e}k\bigg\rceil-1\right).\end{equation}
\end{thm}

\begin{lem}\label{lem: hom} Let $f=\sum_{i=0}^df_ix^i\in\Z[x],\,f_d\neq 0$ be a squarefree polynomial with roots $\alpha_1,\ldots,\alpha_d$ in $\Qb$, $K=\Q(\alpha_1,\ldots,\alpha_d)$ the splitting field of $f$, $p$ a prime number such that $p\nmid f_d\cdot\disc(f)$, $k\ge 1$. Then there exists a ring homomorphism $\phi:\Z[\alpha_1,\ldots,\alpha_d]\to A$, where $A$ is an extension ring of $\Z/p^k\Z$, such that 
\begin{enumerate}
    \item[(i)] $\phi$ induces a bijection between $\{\alpha_1,\ldots,\alpha_d\}$ and the set of roots of $f$ in $A$.
    \item[(ii)] If $\gamma\in\Z[\alpha_1,\ldots,\alpha_d]$ is integral over $\Z$ and $\phi(\gamma)=0$ then $p^k|N_{K/\Q}(\gamma)$.
\end{enumerate}
\end{lem}

\begin{lem}\label{lem: main1} Let $f=\sum_{i=0}^df_ix^i\in\Z[x],\,\deg f=d$ be squarefree and without (rational) integer roots, $S$ the set of roots of $f$ in $\Qb$, $S_1,\ldots,S_r$ a partition of $S$ with $S_i$ potent. There exists a constant $D=D(f)>0$ such that for any sufficiently large $N\ge N_0(f)$ and any prime $p>DN$, $p$ divides at most $d-r$ of the values $f(1),f(2),\ldots,f(N)$.
\end{lem}

\begin{lem}\label{lem: int comb} Let $N,n_1,n_2,\ldots,n_t$ be natural numbers with $n_i\le N$. There exist integers $a_1,\ldots,a_t$, not all zero, such that $|a_i|\le (tN)^{1/t}+1$ and $|a_1n_1+\ldots+a_tn_t|\le (tN)^{1/t}$.\end{lem}

\begin{lem}\label{lem: pre lem main2} Let $f\in\Z[x],\,\deg f=d$ be squarefree and without (rational) integer roots, $S$ the set of roots of $f$ in $\Qb$, $S_1,\ldots,S_r$ a partition of $S$ with $S_i$ potent and $r$ satisfying $\dim\langle S\cup\{1\}\rangle_\Q=d+1-r$. Denote $G_f=\Gal(f/\Q)$.
Then there exists a constant $D=D(f)>0$ such that for any sufficiently large $N> N_0(f)$, a prime $p>DN$ and $k\ge 1$, $p^k$ divides at most
$\big\lceil\frac{|G_f|}k\big\rceil-1$ of the values $f(1),f(2),\ldots,f(N)$.
\end{lem}

\begin{proof} 
Assume that $p^k|f(n_1),\ldots,f(n_t)$ for some $1\le n_i\le N$. We need to show that $t<\frac{|G_f|}k$ whenever $p>DN$ for a suitable constant $D$ and $N$ sufficiently large. Denote by $K/\Q$ the splitting field of $f$. By Lemma \ref{lem: hom}, for $p$ sufficiently large there exists an extension ring $A\supset \Z/p^k\Z$ and a ring homomorphism $\phi:\Z[S]\to A$ with the following properties:
\begin{enumerate}
    \item[(i)] $\phi$ induces a bijection between $S$ and the roots of $f$ in $A$.
    \item[(ii)] For any integral $\gamma\in\Z[S]$ with $\phi(\gamma)=0$ we have $p^k|N_{K/\Q}(\gamma)$.
\end{enumerate}
For any $\alpha\in\Z[S]$ we denote $\bar\alpha=\phi(\alpha)$.
By (i) and the assumption $p^k|f(n_i)$, there exist $\alpha_1,\ldots,\alpha_t\in S$ such that $\bar\alpha_i=\bar n_i$. If we take $D\ge 1$ so that $p>N$, the elements $\alpha_1,\ldots,\alpha_t$ are distinct (since the $\bar n_i$ are). 

 By Lemma \ref{lem: int comb} there exist integers $a_1,\ldots,a_t,\,|a_i|\le (tN)^{1/t}+1$, not all zero, such that $u=\sum_{i=1}^ta_in_i$ satisfies $|u|\le (tN)^{1/t}$. 
Denote by $f_d$ the leading coefficient of $f$ and consider the element
\begin{equation}\label{eq: gamma}\gamma=f_d\sum_{i=1}^ta_i(\alpha_i-n_i)=f_d\sum_{i=1}^t a_i\alpha_i-f_du\in\Z[S].\end{equation}
Since $\bar\alpha_i=\bar n_i$ we have $\bar\gamma=0$. Also $\gamma$ is integral over $\Z$ (because each $f_d\alpha_i$ is) so by property (ii) of $\phi$ we have $p^k|N_{K/\Q}(\gamma).$
 
 We view $\Qb$ as a subset of $\C$, the latter equipped with the usual absolute value. Denote $C=\max_{\alpha\in S}|\alpha|$.
 From (\ref{eq: gamma}) we see that any conjugate $\gamma'$ of $\gamma$ satisfies
 $$|\gamma'|\le f_d\sum_{i=1}^t|a_i|C+f_d|u|\le f_d\cdot\left(((tN)^{1/t}+1)C+(tN)^{1/t}\right)$$ and thus
$$|N_{K/\Q}(\gamma)|\le f_d^{|G_f|}(((tN)^{1/t}+1)C+(tN)^{1/t})^{|G_f|}\le (DN)^{\frac{|G_f|}{t}}$$ if we take $D\ge f_d^t(2C+1)^tt$.
 There are two cases to consider.

\case{$\gamma\neq 0$} In this case we have 
$$0<|N_{K/\Q}(\gamma)|<(DN)^{\frac{|G_f|}t}.$$ 
Now since $p^k|N_{K/\Q}(\gamma)$
and $p>DN$, we must have $t<\frac{|G_f|}k$, proving the assertion of the lemma.

\case{$\gamma=0$} Here we need to use the assumptions (from the statement of Theorem \ref{thm: main2}) that each $S_i$ is potent and 
$\dim\langle S\cup\{1\}\rangle=d+1-r$. Throughout the rest of the proof all spans (denoted by $\langle\cdot\rangle$) and linear dependencies are over $\Q$. 

First we claim that we must have $S_i\subset\{\alpha_1,\ldots,\alpha_t\}$ for some $i$. Otherwise there exists $\beta_i\in S_i\setminus\{\alpha_1,\ldots,\alpha_t\}$ for each $1\le i\le r$ and since $S_i$ is potent $\beta_i\in\langle S_i\setminus\{\beta_i\}\rangle$, so we have
$\langle S\cup\{1\}\rangle=\langle \Cup(S_i\setminus\{\beta_i\})\cup\{1\}\rangle$, with the set in the latter span containing $\alpha_1,\ldots,\alpha_t$. But by (\ref{eq: gamma}) and the assumption $\gamma=0$ we have $\sum_{i=1}^ta_i\alpha_i=u$ and thus $\alpha_1,\ldots,\alpha_t,1$ are linearly dependent over $\Q$, so $$\dim\langle S\cup\{1\}\rangle=\dim\Big\langle \Cup_{i=1}^r(S_i\setminus\{\beta_i\})\cup\{1\}\Big\rangle\le\Bigl|\Cup_{i=1}^r(S_i\setminus\{\beta_i\})\cup\{1\}\Bigr|-1\le d-r,$$ a contradiction to the assumption $\dim\langle S\cup\{1\}\rangle=d+1-r$. This establishes our claim that $S_i\subset\{\alpha_1,\ldots,\alpha_t\}$ for some $i$.

Now suppose (after reordering the $\alpha_j$) that $S_i=\{\alpha_{1},\ldots,\alpha_{{s}}\}$. 
Then since $S_i$ is potent we have $$\sum_{l=1}^{s}b_l\alpha_{l}=w\in\Z$$ for some $0< b_l\in\Z$, and hence $\sum_{l=1}^{s}b_l\bar\alpha_{l}=\bar w$. Since $\bar\alpha_{l}=\bar n_{l}$ we obtain $\sum_{l=1}^{s}b_{l}n_{i}-w\equiv 0\pmod {p}$, which is impossible if $D\ge\sum_{l=1}^{s}b_{l}+1$, $p>DN$ and $N$ is sufficiently large, just as in the proof of Lemma \ref{lem: main1} (here we need to use the assumption that $f$ does not have integer roots). We conclude that the assertion of the lemma holds with $D=\max\left(f_d^t(2C+1)^tt,\sum_{l=1}^sb_l+1\right)$. 


\end{proof}

\end{document}
